%% =====================================================================
%%  T-26-03  No Island
%%  -------------------------------------------------------------------
%%  One idea per file. Core is mandatory. Slots in canonical order.
%%  Max 700 words. No layout commands. No filler. New source material
%%  for this entry goes in sources/B3/T-26-03.tex -- never by
%%  rewriting this file (docs/RULES.md R0.1).
%%  =====================================================================
%% @kind: topic
%% @id: T-26-03
%% @title: No Island
%% @subtitle: Avoid the unhappy and unlucky --- misery is infectious, and its carriers are not liars
%% @chapter: c26
%% @order: 30
%% @status: stable
%% @sources: B3
%% @updated: 2026-09-19

\topic{T-26-03}{No Island}{Avoid the unhappy and unlucky --- misery is infectious, and its carriers are not liars}

\begin{core}
You can die of someone else's misery: emotional states are as infectious as diseases, and the drowning person you are trying to save may be the weight that drowns you. B3's cold instruction is to associate with the happy and fortunate instead --- and the exit discipline that follows is genuinely hard, because the infectious are not faking, their misfortunes are often real, and pity is exactly how the infection spreads.
\end{core}
\begin{mechanism}
The source's law of infection runs on proximity: the associate with the turbulent character --- Lola Montez is the type, leaving a trail of ruined careers and short tempers --- ends up paying for the turbulence as if it were their own, because shared hours become shared weather. Two mechanisms drive it. Contagion: moods and causal entanglements both transfer --- the unlucky carry their storm systems into every room and every ledger they touch, and their crises become your errands. Enfeeblement: chronic exposure to the unhappy recalibrates your own baseline --- what you tolerate rises, what you expect falls, and the association's costs are paid in drift before they are paid in disasters. The source adds the diagnostic trap: the infectious often present as victims, which makes their miseries hard to read as self-inflicted until you are already entangled --- the rescue fantasy is the disease's preferred vector. The reversal is narrow and honest: this law is about chronic infection, not about abandoning people in a hard season; the test is trajectory and pattern, not suffering itself.
\end{mechanism}
\begin{conditions}
Strongest for close, sustained associations --- partners, co-founders, housemates, the daily colleague --- where contagion compounds. It weakens at distance and with bounded, professional contact, and it inverts where the unhappiness is situational and being actively repaired: the law prices patterns, not moments. It is at its most dangerous when it is at its most sympathetic.
\end{conditions}
\begin{application}
  \item Audit the five people you spend most hours with: how many are in a storm, and is the storm theirs to end or a climate?
  \item Distinguish the season from the pattern: a person in a bad year with repairs underway is a neighbour; a person whose every year is someone else's fault is a climate.
  \item Help by design, not by immersion: bounded help with a defined end (money lent once, advice given on request) does not transfer the weather.
  \item When you must exit, exit without a trial. The infected experience departure as abandonment; a prolonged goodbye becomes the next entanglement.
  \item Watch your own climate: the unhappy often became so by association. Your circle is your forecast.
\end{application}
\begin{feedback}
  \item Working: your baseline recovers --- plans feel makeable again within a season of the exit.
  \item Working: your bounded help actually helped, and the boundary held. The rare good outcome of this entry.
  \item Failing: you are running on their crisis schedule, and your own projects have a taste of adrenaline and dread.
  \item Failing: you notice your other friendships thinning --- the storm system crowds out the clear skies first.
\end{feedback}
\begin{failure}
The doctrine's abuse is its known use: rules about avoiding the unlucky have always been the machinery of abandonment, and the vocabulary of infection makes cruelty sound like hygiene. The practitioner's real discipline is holding both truths --- misery is contagious, and the sick still deserve care --- which means bounded help as the default, exit as the exception, and the honest admission that this entry will be cited by people who just did not want the inconvenience.
\end{failure}
\begin{countermeasures}
  \item Notice who leaves your presence smaller. The bodily signal is honest weeks before the ledger is.
  \item Track help you have given: if it converts repeatedly into entitlement and deeper crisis, you are not helping a person, you are feeding a pattern (T-19-02).
  \item Refuse the rescue fantasy specifically: \emph{only I can save them} is the infection speaking in your own voice, wearing your generosity.
  \item If you are the storm, the countermeasure is professional repair and a long apology to the people whose weather you made. The law applies to you from the inside too.
\end{countermeasures}

\sources{B3}
\seesources
\seealso{T-03-01, T-27-03, T-26-02}
